
Research automation's validation crisis was addressed: automated ML pipelines fail late when constraint violations propagate undetected through phase boundaries. Contract-based validation addresses this by enforcing semantic and compositional constraints at phase boundaries, preventing violations from reaching downstream implementation phases.

The key insight---that three-layer validation (schema + pattern + compositional contracts) forces explicit checking of constraints schema-only validation cannot express---led to a practical solution. Formal Design-by-Contract methods, traditionally applied to code correctness, were demonstrated to transfer to research workflow automation when outputs are structured and constraints are explicit. The result: 100\% downstream failure reduction on a 100-case placeholder hypothesis corpus with four feasibility constraints on placeholder content with typed interfaces (generalization to substantive research unproven, external validity deferred to future work).

The causal mechanism holds across all validation steps. Contracts force explicit checking (100 percentage point expressiveness gap over schema-only). Multi-layer architecture achieves complementary detection (88\% versus 40\%, with cumulative layer contributions). Early boundary validation prevents propagation (100\% of violations caught before Phase 4/5 implementation). Each mechanism step was validated independently before testing end-to-end effectiveness.

Three findings merit emphasis. First, feasibility constraints are fundamentally semantic and compositional, not structural---schema validation checking types and presence cannot express keyword patterns, cross-field logic, or state-based membership. Second, validation layers must be complementary not redundant, with each layer targeting violations the previous layer cannot detect. Third, fail-fast at boundaries changes the failure mode from expensive late-stage discovery (Phase 4/5 implementation failures) to immediate early rejection (Phase $2\rightarrow 3\rightarrow 4$ boundary validation), shifting the cost from hours of wasted compute to milliseconds of checking.

The limitations are clear. Testing was on placeholder content with minimal research fidelity, demonstrating that constraint enforcement logic works on typed interfaces but not proving generalization to substantive research with complex semantic dependencies. Pattern layer recall (80\% on keyword constraints) falls short of the 95\% target due to synonym gaps in finite keyword blacklists. Manual contract specification (7 LOC per constraint) does not scale to workflows with 10+ constraints without automation. External validity to production ML pipelines with real research content remains unproven.

Future work opens multiple directions. Auto-generation of contracts from natural language constraint descriptions would eliminate manual specification burden, enabling contract-based validation for large constraint sets. LLM-based semantic validation layers could replace pattern matching with prompted classification, achieving 95\%+ recall by handling paraphrases and synonyms pattern matching misses. Cross-phase compositional contracts maintaining state across multiple transitions could prevent constraint drift (Phase 2 specifies CIFAR-10, Phase 4 accidentally uses MNIST). Extension to model training pipelines---validating data split disjointness, hyperparameter bounds, architecture constraints---would test whether contract-based validation generalizes beyond research automation to production ML workflows.

The path from formal methods for code to formal methods for research automation is now open. Three-layer validation demonstrates that Design-by-Contract principles enable automated research pipelines to enforce constraints without human review on placeholder content with typed interfaces, achieving zero downstream failures on infrastructure testing (external validity to substantive research requires future validation). From constraint violation crisis to zero downstream failures on placeholder workflows---that is the contribution of contract-based validation.
